\documentclass[10pt,a4paper]{amsart}
\usepackage[margin=19mm]{geometry}
\usepackage{amsmath,amssymb,mathtools}
\usepackage[scr=rsfs,cal=euler]{mathalfa}
\usepackage{xfrac}
\usepackage[unicode,hidelinks,bookmarks=false]{hyperref}
\newcommand{\ie}{i.e.\ }
\newcommand{\cf}{cf.\ }
\newcommand{\resp}{respectively\ }
\newcommand{\Subsection}{\subsection}
\providecommand{\nobreakdash}{\nobreak\hbox{-}}
\setlength{\emergencystretch}{2em}
\multlinegap0pt
\usepackage{bm}
\usepackage{bbm}
\usepackage{dsfont}
\usepackage{xspace}
\usepackage{cancel}
\usepackage{mathtools}
\usepackage{amsthm}
\makeatletter
\@ifundefined{result}{\newtheorem{result}{Result}}{}
\@ifundefined{lem}{\newtheorem{lem}[result]{Lemma}}{}
\makeatother
\makeatletter
\expandafter\ifx\csname proof\endcsname\relax
\renewenvironment{proof}[1][\proofname]{\par
  \vspace{-\topsep}% remove the space after the theorem
  \pushQED{\qed}%
  \normalfont
  \topsep0pt \partopsep0pt % no space before
  \trivlist
  \item[\hskip\labelsep
        \itshape
    #1\@addpunct{.}]\ignorespaces
}{%
  \popQED\endtrivlist\@endpefalse
  \addvspace{6pt plus 6pt} % some space after
}
\fi
\makeatother
\title[Three-color van der Waerden numbers grow super-exponentially]{Three-color van der Waerden numbers grow super-exponentially}
\author{Jacob Fox, Zach Hunter}
\date{Source: arXiv:2606.02541v1 (CC BY 4.0)}
\begin{document}
\maketitle

\noindent\emph{Source and licence.} Three-color van der Waerden numbers grow super-exponentially. Version arXiv:2606.02541v1 (CC BY 4.0), distributed under the Creative Commons Attribution 4.0 International licence, as recorded in the arXiv licence metadata for this exact version and shown on the abstract page https://arxiv.org/abs/2606.02541v1. Full rights notice: licenses/arxiv-2606.02541-CC-BY-4.0.txt.\par\smallskip

\noindent\textit{Editorial scope.} Counted unit: the Lem whose source label is \texttt{intervaldense} in the article (source0.tex lines 354--356, source label \texttt{intervaldense}), delivered together with the article's own complete original proof of it (source0.tex lines 357--376, one \texttt{proof} environment of 20 lines). No mathematics is written, repaired or re-derived: statement and proof are the article's own text line for line. Required context retained uncounted: none. Every definition, display and label of those lines is retained unchanged. Omitted from this package and not redistributed: 1--353, 377--745. Nothing inside a retained excerpt is omitted or altered: every retained body is delivered exactly as the article's lines are written, so no omission receipt of any kind accompanies this package. Citations: the retained excerpts carry no citation command at all, so no bibliography entry of the article is transcribed and no reference is redistributed. Reference adaptation: none was needed and none was made; every reference inside the delivered unit resolves inside it, so no unresolved reference or citation is produced. Printed slips kept literal: no printed word, symbol, hypothesis or number of the counted statement or of its proof was altered. Layout and class: the article's own class is \texttt{article} with packages including inputenc, amsmath, amssymb, amsfonts, amsthm, mathrsfs, bm, enumerate, bbm, tikz, csvsimple, centernot, hyperref, algorithm2e; the wrapper is the lane's pinned \texttt{amsart} editorial wrapper at 10pt on A4, which restates the theorem-like environments the retained text needs and adds the article's own macro definitions that the retained text needs (none), with extra wrapper packages (bm, bbm, dsfont, xspace, cancel, mathtools). The delivered numbering is the unit's own. Assets: no retained span contains a figure environment, a graphics-inclusion command, a PDF-page inclusion, a table or a TikZ picture; no image file is copied into this package, the package declares no asset dependency and no image file is redistributed with it. Licence: version arXiv:2606.02541v1 of this article is distributed under the Creative Commons Attribution 4.0 International licence, as recorded in the arXiv licence metadata for this exact version and shown on the abstract page https://arxiv.org/abs/2606.02541v1. A full rights notice is delivered at licenses/arxiv-2606.02541-CC-BY-4.0.txt. Characters: every retained excerpt is pure ASCII, so the delivered engine, XeLaTeX with its default Latin Modern text font, reports no missing character.
\begin{lem}\label{intervaldense}
For each sufficiently small $\varepsilon>0$ there is $\delta>0$ such that the following holds. For any integer $k \geq e^{36/\varepsilon}$ and integer $k \leq N \leq 2.5^k$, there is a subset $S \subseteq [N]$ with $|S| \leq \varepsilon N$ that contains at least a $\delta$-fraction of any $k$-term arithmetic progression in $[N]$. Furthermore, we can take $\delta =2^{-400/\varepsilon}$. 
\end{lem}

\begin{proof}
Let $\ell=\lceil k/101 \rceil$ and $Q$ denote the set of primes in the interval $(\ell,k]$. By the prime number theorem, $|Q| \geq 0.99k/\log k$. Let $t=\lfloor \varepsilon (\log k)/6\rfloor \geq \varepsilon (\log k)/7$, where the inequality uses $k \geq e^{36/\varepsilon}$. 

Independently for each prime $p \in Q$, choose $t$ distinct congruence classes modulo $p$ uniformly at random, and let $S_p$ be the elements in $[N]$ that are in at least one of these $t$ congruence classes. Let $S=\bigcup_{p \in Q} S_p$. 

{\bf Step 1: Density of $S$.} We first show that $|S| \leq \varepsilon N$. Indeed, $$|S|\le \sum_{p\in Q} |S_p|\le \sum_{p\in Q} (N/p+1)t = |Q|t+tN \sum_{p\in Q}1/p\le \varepsilon k /6+4.7\varepsilon N/6 < \varepsilon N.$$
To see the second to last inequality, let $$f(x)=\sum_{\substack{p~\textrm{prime} \\ p \leq x}} 1/p,$$ so $$\sum_{p \in Q} 1/p=f(k)-f(\ell)=\frac{(1+o(1))}{\log k}\sum_{n=k/101}^k 1/n = (1+o(1))\frac{\log(101)}{\log k}<4.7/\log k,$$ where the second equality holds by the prime number theorem, and the final inequality is due to $k$ being sufficiently large (which in turn holds as $\varepsilon$ is assumed to be sufficiently small).



{\bf Step 2: Discrepancy of $S$ in $k$-APs.} We next show that with positive probability the set $S$ has density at least $\delta$ in any $k$-term arithmetic progression in $[N]$. The proof is by a union bound. Fix a $k$-term arithmetic progression $P$ with elements in $[N]$ and let $d$ denote the common difference of $P$. Fix a subset $D$ of the arithmetic progression $P$ with size $u=\lfloor \delta k \rfloor$. 

If $d$ is a multiple of $s$ primes in $Q$, then the product of these primes is a factor of $d$ and hence $d \geq \ell^s$. As $N > d$, we have $s \leq \log_{\ell} N \leq k(\log 2.5)/\log \ell$. Let $Q'$ be the subset of $Q$ of primes $p$ that are not a factor of $d$. So $$|Q'| \geq 0.99k/\log k - k(\log 2.5)/\log \ell \geq k/(14\log k),$$using $\log\ell\ge\log k-\log(101)$ and the largeness of $k$. For each prime $p \in Q'$, we partition $P$ into subsets $P_i$ with $i \in \mathbb{Z}_p$ being those elements of $P$ that are congruent to $i \pmod p$. 

Each $P_i$ has size at least $\lfloor k/p \rfloor \geq k/(2p)$ as $p \leq k$. So the set $D$, of size $\delta k$, can contain all elements of at most $2\delta p$ of the subsets $P_i$. So the probability that $S_p \cap P \subseteq D$ is at most $(2\delta)^t$. As these choices are done independently for each prime $p \in Q'$, the probability that $S \cap P \subseteq D$ is at most $((2\delta)^t)^{|Q'|} \leq (2\delta)^{\varepsilon k/98}$. 

The number of choices of the progression $P$ is at most $N^2$. The number of choices for the subset $D$ of $P$ is ${k \choose u} \leq 2^k$. Hence, by the union bound, the probability that $S$ has density less than $\delta$ in some $k$-term arithmetic progression in $[N]$ is at most $$N^2 2^k (2\delta)^{\varepsilon k/98}<(2.5)^{2k}2^k(2^{1-400/\varepsilon})^{\varepsilon k/98}<1.$$

Since $S$ has the desired properties with positive probability, there exists such a desired set $S$, which completes the proof. 
\end{proof}

\end{document}
